% !TeX root = ../../main.tex

\section{Limitations}
\label{sec:limitations}

It should be noted that VisR and DoPIo are similar in their use of on-axis ARF-induced displacements for assessing tissue mechanical properties. As such, both techniques, as well as the experiments that characterized and compared them, share several overarching limitations. These limitations can be broadly categorized into those related to beam sequences, assumptions about material properties (i.e. the tissue or other media being interrogated), and experimental design.

\subsection{Beam Sequence Limitations}

Regarding beam sequences, both VisR and DoPIo rely on the use of ARF excitations to induce on-axis displacements in tissue. One major variable for how ARF pushes are generated, as well as how they are tracked, is the choice of beam widths (i.e. F-numbers) used for the push and track beams. The experiments in this dissertation used F/1.5 pushes in VisR, as described in previous works~\cite{hossain_Evaluating_2018,hossain_Viscoelastic_2020}, whereas DoPIo was found to perform best with an F/3.0 push. The manner in which this work understood these choices primarily relied on the idea that the initial displacements that are tracked shortly after the ARF excitation are governed by the PSF of the push beam. This interpretation ultimately arises from the understanding that displacements are induced by ARF by way of mechanical energy transfer from the acoustic beam to the medium~\cite{sarvazyan_Shear_1998,sarvazyan_Acoustic_2021}. However, this is only one of four mechanisms in which ARF may be generated in tissue-like media; in \cite{sarvazyan_Acoustic_2021}, Sarvazyan \emph{et al.} describes this mechanism as "A-ARF". The other three mechanisms are "B-ARF", which arise from the reflection of acoustic waves at interfaces (as seen in radiation force balances); "C-ARF", which arise from spatial variations in energy density from standing waves (e.g. in acoustic tweezers); and "D-ARF", which arise from variations in energy density due to gradients in mass density or compressional wave speeds. Since the renal parenchyma is a complex, heterogeneous medium, it is possible that B-ARF and D-ARF mechanisms may also contribute to the overall displacement fields that were used to calculate VisR MSD parameters and DoPIo times-intersect (and, therefore, RE/RV parameters for VisR and modulus estimates for DoPIo). The significance of these additional ARF mechanisms on displacement profiles for generic ARFI acquisitions is not well understood, let alone their specific effects on VisR and DoPIo measurements. Further investigation is warranted to elucidate the relevance of these other two ARF mechanisms as confounding factors, as well as their relevance specifically to VisR and DoPIo measurements.

Another limitation is that both VisR and DoPIo rely on the assumption that the ARF-induced displacements being tracked are purely axial. Although ARF excitations actually induce complex 3D displacement fields in tissue, it is commonly assumed that displacements in directions orthogonal to the ARF push beam are at least an order of magnitude smaller such that they may be neglected for ARFI-based applications~\cite{palmeri_Experimental_2004,palmeri_Ultrasonic_2006}. However, this assumption is primarily motivated by an application to off-axis, shear wave-based techniques; the insignificance of lateral and elevational displacements to on-axis techniques like VisR and DoPIo has not been confirmed in the literature. Future work should investigate the presence of off-axis displacements in ARFI sequences used for VisR and DoPIo, as well as their potential effects on measurement accuracy.

Likewise, the PSF of push and track beams are assumed to be invariant across space (i.e. across different A-lines), time (i.e. across different acquisition instances), subjects (i.e. different patients or animals), and imaging plane orientations (i.e. 0 versus 90 degrees). However, in practice, these factors may all influence the effective PSF of the beams used for ARF excitation and displacement tracking. For example, variations in tissue attenuation and near-field geometries (e.g. skin, fat, muscle layers) across subjects and imaging planes may distort acoustic energy distributions over space, thereby changing the effective PSF of both push and track beams. These variations could introduce inconsistencies in VisR and DoPIo measurements that are not accounted for in the current implementations. Future studies should consider characterizing and compensating for these potential sources of variability to enhance measurement reliability.

\subsection{Material Property Assumptions}

Both VisR and DoPIo rely on specific assumptions about the material properties of the media being interrogated. VisR assumes that the medium behaves as a linear, viscoelastic material that can be modeled using a one-dimensional MSD model, whereas DoPIo is developed using FEM simulations of a linear, isotropically elastic medium. However, as mentioned in the introduction, biological tissues such as renal parenchyma exhibit complex, nonlinear, anisotropic, and viscoelastic behaviors that may not be fully captured by these simplified models. It should be noted that this dissertation sought to control for mechanical anisotropy by controlling transducer angle orientations during data acquisition. In addition, viscoelastic effects were partially addressed by derived RE and RV parameters in VisR.\footnote{
    A small but growing body of work has begun to investigate the effects of nonlinear elasticity using ARF-based elastography techniques. However, further research is needed to understand the clinical relevance of those parameters for nephropathological applications.
} 

While it can be helpful to characterize the limitations of VisR and DoPIo in the context of their underlying material property assumptions, it is equally concievable that these techniques may be augmented such that they are specifically geared towards use in more complex media. Since VisR lookup tables for elasticity correction and DoPIo training data for empirical models have both been developed purely in isotropic media, future work may explore modifications to those frameworks that incorporate anisotropic and/or nonlinear material properties. This may involve generating new FEM simulation datasets that account for these complexities, then using those datasets to retrain DoPIo empirical models or regenerate VisR lookup tables. Such efforts could potentially enhance the accuracy and clinical utility of both techniques when applied to biological tissues like renal parenchyma.

\subsection{Experimental Design Limitations}

Finally, both VisR and DoPIo experiments in this dissertation share several limitations related to experimental design. For instance, the sample sizes for patients with clinically indicated biopsies in Chapter~\ref{chap:visr} (particularly for those with rejection episodes) and pigs in Chapter~\ref{chap:iri} were relatively small. This limited sample size impacts the ability to generalize the results with sufficient statistical power. Furthermore, in order to make those measurements in the first place, a significant number of FEM simulations were required to develop empirical models for DoPIo and elasticity correction lookup tables for VisR. As a result, existing models are not valid outside the specific ranges of material properties and experimental conditions that were simulated. This manifests as limits to the range and resolution of elastic moduli that may be estimated using VisR and DoPIo (and also viscosity for VisR), as well as an inability for DoPIo to distinguish between stiffer, more viscous materials versus softer, less viscous materials that may produce similar displacement profiles.

The most obvious solution to these limitations involves generating larger datasets with more diverse material properties and experimental conditions. However, this is easier said than done. The computational resources and time required to generate these simulations limits the accessibility of these techniques for broader research and clinical applications, as well as the ability to explore a wider range of material properties and experimental conditions (especially different configurations for ARF push and track beams, including F-numbers, frequencies, beamforming algorithms, and even transducer geometries). Future work should focus on optimizing experimental designs to increase sample sizes, improve statistical power, and reduce computational burdens associated with model development.

Machine learning techniques may provide a promising avenue for addressing these limitations regarding empirical model scope and generalizability. For example, displacement profiles could be directly inputted into deep learning architectures (e.g. convolutional neural networks, recurrent neural networks) to predict shear elastic and viscous moduli without relying on simplified material property assumptions or extensive FEM simulations. Proofs-of-concept have been developed for VisR using fully-connected neural networks by Richardson \emph{et al.}~\cite{richardson_Quantitative_2024}, as well as for DoPIo using a support vector regressor by Rahmouni \emph{et al.}~\cite{rahmouni_Machine_2020}. Considering that even relatively shallow and simple models have demonstrated superior performance over the traditional MSD- and time-intersect-based approaches for VisR and DoPIo respectively, future work may explore more complex architectures and larger training datasets to further enhance measurement accuracy and generalizability. Depending on the model architecture and input features used, one may conceive of developing a unified model where single ARFI displacement profiles may be used to predict both shear elastic and viscous moduli simultaneously. Alternatively, models may be trained using a diverse enough dataset such that they may generalize across a wide range of ARF push and track beam configurations - or perhaps, even in cases where alternative sources of mechanical excitation (e.g. external mechanical vibrators, endogenous physiological motion) are used to induce displacements. These directions may help overcome some of the limitations associated with the current implementations of VisR and DoPIo, thereby enhancing their clinical utility for renal allograft assessment.
